\documentclass[11pt]{article}

\usepackage[margin=1in]{geometry}
\usepackage{enumitem}
\usepackage{booktabs}
\usepackage{longtable}
\usepackage{array}
\usepackage{amsmath,amssymb}
\usepackage{hyperref}
\usepackage{xcolor}

\title{\textbf{Calculus I for Biological and Social Sciences}\\
Course Structure and Full Topic Breakdown}
\author{}
\date{}

\begin{document}

\maketitle

\section{Course Description}

This course is an applied Calculus I course designed for students in biological and social sciences. It develops the mathematical language and tools used to analyze change, growth, accumulation, optimization, and quantitative models.

The course begins with essential precalculus review and functions, then develops limits, derivatives, and integrals. It emphasizes interpretation, graphical reasoning, modeling, and applications such as population growth, drug absorption, biological systems, rates of change, and cumulative change.

Although the textbook includes later material on differential equations, linear algebra, multivariable calculus, systems, probability, and statistics, the core Calculus I portion is primarily contained in Chapters 1--7, with selected introductory applications from Chapter 8 where appropriate.

\section{Course Structure at a Glance}

\begin{table}[h!]
\centering
\begin{tabular}{p{0.12\textwidth} p{0.27\textwidth} p{0.52\textwidth}}
\toprule
\textbf{Module} & \textbf{Main Theme} & \textbf{Primary Mathematical Ideas} \\
\midrule
Module 0 & Mathematical Foundations & Algebra, graphs, functions, exponentials, logarithms, trigonometry \\
Module 1 & Discrete Models and Sequences & Exponential growth, recurrence equations, sequences, discrete logistic models \\
Module 2 & Limits and Continuity & Limits, continuity, limits at infinity, trigonometric limits \\
Module 3 & Differential Calculus & Derivative definition, rules of differentiation, chain rule, implicit differentiation \\
Module 4 & Applications of Derivatives & Optimization, graph analysis, extrema, concavity, related rates, approximation \\
Module 5 & Integral Calculus & Definite and indefinite integrals, Fundamental Theorem of Calculus, accumulation \\
Module 6 & Integration Techniques and Numerical Methods & Substitution, integration by parts, partial fractions, numerical integration \\
Module 7 & Introductory Differential Equations and Modeling & Separable differential equations, equilibria, stability, population and compartment models \\
\bottomrule
\end{tabular}
\end{table}

\section{Major Course Learning Outcomes}

By the end of the course, students should be able to:

\begin{enumerate}[label=\arabic*.]
    \item Use algebraic, graphical, numerical, and verbal representations to describe functions and mathematical models.
    \item Analyze polynomial, rational, exponential, logarithmic, and trigonometric functions.
    \item Construct and interpret discrete-time models, sequences, and recurrence equations.
    \item Evaluate limits and determine continuity of functions.
    \item Interpret derivatives as instantaneous rates of change and slopes of tangent lines.
    \item Differentiate algebraic, exponential, logarithmic, trigonometric, and implicitly defined functions.
    \item Apply derivatives to optimization, related-rate problems, curve sketching, and approximation.
    \item Interpret and compute definite integrals as accumulated change and area.
    \item Use the Fundamental Theorem of Calculus to connect derivatives and integrals.
    \item Apply integration techniques to solve applied problems in biological and social sciences.
    \item Use numerical methods when exact symbolic solutions are impractical.
    \item Formulate and interpret simple differential-equation models, including growth, decay, compartment, and population models.
\end{enumerate}

\section{Prerequisite Knowledge}

\subsection{Required Mathematical Prerequisites}

Students should be comfortable with the following topics before beginning the course:

\begin{itemize}
    \item Arithmetic with fractions, decimals, percentages, and signed numbers.
    \item Solving linear equations and inequalities.
    \item Manipulating algebraic expressions.
    \item Factoring polynomials.
    \item Working with powers, roots, and rational exponents.
    \item Solving quadratic equations by factoring, completing the square, and the quadratic formula.
    \item Understanding coordinate planes, slope, intercepts, and equations of lines.
    \item Graphing basic functions.
    \item Basic trigonometric functions, especially sine, cosine, tangent, and radians.
    \item Exponential and logarithmic expressions.
    \item Function notation, domain, range, and inverse functions.
\end{itemize}

\subsection{Recommended Scientific and Applied Skills}

Because this course is intended for biological and social science students, the following skills are especially useful:

\begin{itemize}
    \item Interpreting tables, graphs, and real-world data.
    \item Working with units and unit conversions.
    \item Explaining results in context.
    \item Using a calculator, graphing calculator, spreadsheet, or mathematical software.
    \item Reading simple quantitative models.
    \item Understanding common scientific variables such as population, concentration, time, rate, and total accumulation.
\end{itemize}

\section{Detailed Module Breakdown}

\subsection{Module 0: Mathematical Foundations and Functions}

\textbf{Textbook basis:} Chapter 1, ``Preview and Review.''

\subsubsection*{Major Topics}

\begin{itemize}
    \item Real numbers and interval notation.
    \item Lines in the plane.
    \item Equations of circles.
    \item Trigonometry.
    \item Exponential and logarithmic functions.
    \item Complex numbers and quadratic equations.
    \item Elementary functions.
    \item Polynomial, rational, power, exponential, logarithmic, and trigonometric functions.
    \item Inverse functions.
    \item Function transformations.
    \item Graphing and logarithmic scales.
    \item Transforming nonlinear relationships into linear forms.
\end{itemize}

\subsubsection*{Core Concepts}

\begin{itemize}
    \item A function assigns one output to each allowed input.
    \item The domain determines which inputs are meaningful.
    \item Graphs communicate relationships between variables.
    \item Exponential functions model constant proportional change.
    \item Logarithms undo exponential functions and are useful for analyzing multiplicative relationships.
    \item Transformations shift, stretch, reflect, or otherwise modify graphs.
    \item Logarithmic scales help visualize data spanning several orders of magnitude.
\end{itemize}

\subsubsection*{Learning Objectives}

Students will be able to:

\begin{enumerate}[label=\arabic*.]
    \item Determine domains and ranges of elementary functions.
    \item Evaluate functions using equations, tables, and graphs.
    \item Identify intercepts, slopes, asymptotes, and key graph features.
    \item Solve linear and quadratic equations.
    \item Use exponential and logarithmic laws correctly.
    \item Graph and transform elementary functions.
    \item Recognize when a logarithmic transformation can linearize data.
    \item Interpret functions in applied contexts.
\end{enumerate}

\subsubsection*{Potentially Difficult Topics}

\begin{itemize}
    \item Domain restrictions for rational, logarithmic, and root functions.
    \item Inverse functions and one-to-one behavior.
    \item Distinguishing exponential growth from linear growth.
    \item Applying logarithm laws correctly.
    \item Understanding graph transformations.
    \item Using logarithmic scales and transformed variables.
\end{itemize}

\subsection{Module 1: Discrete-Time Models, Sequences, and Difference Equations}

\textbf{Textbook basis:} Chapter 2.

\subsubsection*{Major Topics}

\begin{itemize}
    \item Exponential growth and decay.
    \item Population growth in discrete time.
    \item Recurrence equations.
    \item Sequences.
    \item Recursive sequences.
    \item Limits of sequences.
    \item Summation notation.
    \item Density-dependent population growth.
    \item Beverton--Holt population model.
    \item Discrete logistic equation.
    \item Applied models such as drug absorption.
\end{itemize}

\subsubsection*{Core Concepts}

A discrete-time model describes quantities measured at separate time intervals, such as daily population counts, annual economic data, or weekly disease cases.

A general recurrence equation has the form:
\[
x_{n+1} = f(x_n),
\]
where \(x_n\) is the quantity at time step \(n\).

Exponential growth in discrete time is commonly represented by:
\[
P_{n+1} = rP_n,
\]
or equivalently,
\[
P_n = P_0r^n.
\]

A discrete logistic model includes density dependence:
\[
P_{n+1} = rP_n\left(1-\frac{P_n}{K}\right),
\]
where \(K\) represents a carrying-capacity-like limiting quantity.

\subsubsection*{Learning Objectives}

Students will be able to:

\begin{enumerate}[label=\arabic*.]
    \item Construct recurrence equations from verbal descriptions.
    \item Distinguish between linear, exponential, and density-dependent growth.
    \item Calculate terms of recursive sequences.
    \item Use summation notation to represent accumulated quantities.
    \item Interpret long-term behavior of simple discrete models.
    \item Model population change and other repeated processes.
    \item Compare unrestricted growth with limited growth models.
\end{enumerate}

\subsubsection*{Potentially Difficult Topics}

\begin{itemize}
    \item Translating a word problem into a recurrence equation.
    \item Understanding the difference between discrete and continuous growth.
    \item Interpreting parameters in population models.
    \item Finding and interpreting sequence limits.
    \item Understanding density dependence and carrying capacity.
    \item Distinguishing recursive formulas from explicit formulas.
\end{itemize}

\subsection{Module 2: Limits and Continuity}

\textbf{Textbook basis:} Chapter 3.

\subsubsection*{Major Topics}

\begin{itemize}
    \item Informal meaning of limits.
    \item Algebraic evaluation of limits.
    \item Limit laws.
    \item One-sided limits.
    \item Limits involving infinity.
    \item Continuity.
    \item Intermediate Value Theorem.
    \item Bisection method.
    \item Trigonometric limits.
    \item Sandwich Theorem.
    \item Formal definition of a limit.
\end{itemize}

\subsubsection*{Core Concepts}

The notation
\[
\lim_{x\to a}f(x)=L
\]
means that \(f(x)\) approaches \(L\) as \(x\) approaches \(a\), even if \(f(a)\) is undefined or differs from \(L\).

A function \(f\) is continuous at \(x=a\) when:
\[
\lim_{x\to a}f(x)=f(a).
\]

Continuity is important because it supports major existence results, including the Intermediate Value Theorem. If a continuous function changes from positive to negative on an interval, then it must equal zero somewhere in that interval.

\subsubsection*{Learning Objectives}

Students will be able to:

\begin{enumerate}[label=\arabic*.]
    \item Explain the difference between a function value and a limit.
    \item Estimate limits from graphs, tables, and numerical evidence.
    \item Evaluate limits algebraically using limit laws.
    \item Identify removable, jump, and infinite discontinuities.
    \item Determine continuity at a point and on an interval.
    \item Evaluate limits at infinity.
    \item Apply the Intermediate Value Theorem.
    \item Use the bisection method to approximate solutions of equations.
\end{enumerate}

\subsubsection*{Potentially Difficult Topics}

\begin{itemize}
    \item Understanding that a limit can exist when a function is undefined.
    \item Evaluating indeterminate forms such as \(0/0\).
    \item Distinguishing vertical asymptotes from removable discontinuities.
    \item Limits at infinity and horizontal asymptotes.
    \item Trigonometric limits, particularly:
    \[
    \lim_{x\to 0}\frac{\sin x}{x}=1.
    \]
    \item Formal limit definitions using \(\varepsilon\) and \(\delta\).
\end{itemize}

\subsection{Module 3: Differential Calculus}

\textbf{Textbook basis:} Chapter 4.

\subsubsection*{Major Topics}

\begin{itemize}
    \item Formal definition of the derivative.
    \item Derivative as slope and instantaneous rate of change.
    \item Differentiability and continuity.
    \item Power rule.
    \item Product rule.
    \item Quotient rule.
    \item Chain rule.
    \item Implicit differentiation.
    \item Related rates.
    \item Higher derivatives.
    \item Derivatives of trigonometric functions.
    \item Derivatives of exponential and logarithmic functions.
    \item Derivatives of inverse functions.
    \item Linear approximation and error propagation.
\end{itemize}

\subsubsection*{Core Concepts}

The derivative of \(f\) at \(x\) is defined by:
\[
f'(x)=\lim_{h\to 0}\frac{f(x+h)-f(x)}{h}.
\]

The derivative measures instantaneous change. Depending on context, it can represent:

\begin{itemize}
    \item Velocity as the derivative of position.
    \item Population growth rate as the derivative of population.
    \item Marginal cost as the derivative of cost.
    \item Drug concentration change as the derivative of concentration.
    \item The slope of the tangent line to a graph.
\end{itemize}

Important differentiation rules include:
\[
\frac{d}{dx}\left(x^n\right)=nx^{n-1},
\]
\[
\frac{d}{dx}\left(f(x)g(x)\right)=f'(x)g(x)+f(x)g'(x),
\]
\[
\frac{d}{dx}\left(\frac{f(x)}{g(x)}\right)
=
\frac{f'(x)g(x)-f(x)g'(x)}{[g(x)]^2},
\]
and
\[
\frac{d}{dx}f(g(x))=f'(g(x))g'(x).
\]

\subsubsection*{Learning Objectives}

Students will be able to:

\begin{enumerate}[label=\arabic*.]
    \item Compute derivatives from the limit definition.
    \item Interpret derivatives graphically and in context.
    \item Determine whether a function is differentiable.
    \item Differentiate polynomial, rational, power, exponential, logarithmic, and trigonometric functions.
    \item Apply the product, quotient, and chain rules.
    \item Differentiate implicitly defined functions.
    \item Solve related-rate problems with appropriate units.
    \item Calculate higher derivatives.
    \item Construct tangent-line approximations.
    \item Estimate propagated error using linearization.
\end{enumerate}

\subsubsection*{Potentially Difficult Topics}

\begin{itemize}
    \item Computing derivatives directly from the limit definition.
    \item Recognizing when to use the product, quotient, or chain rule.
    \item Applying the chain rule multiple times in nested functions.
    \item Implicit differentiation.
    \item Related rates, especially identifying variables that change over time.
    \item Derivatives involving logarithms, exponentials, and inverse functions.
    \item Distinguishing \(dy/dx\), \(\Delta y\), and differentials.
\end{itemize}

\subsection{Module 4: Applications of Differentiation}

\textbf{Textbook basis:} Chapter 5.

\subsubsection*{Major Topics}

\begin{itemize}
    \item Absolute and local extrema.
    \item Extreme Value Theorem.
    \item Mean Value Theorem.
    \item Increasing and decreasing functions.
    \item Concavity.
    \item Inflection points.
    \item Optimization.
    \item L'H\^{o}pital's Rule.
    \item Curve sketching.
    \item Asymptotes.
    \item Newton--Raphson method.
    \item Stability of recurrence equations.
    \item Antiderivatives.
    \item Applications to biological systems and drug models.
\end{itemize}

\subsubsection*{Core Concepts}

Critical points occur where:
\[
f'(x)=0
\]
or where \(f'(x)\) does not exist.

The first derivative provides information about increasing and decreasing behavior:

\[
f'(x)>0 \Rightarrow f \text{ is increasing},
\]
\[
f'(x)<0 \Rightarrow f \text{ is decreasing}.
\]

The second derivative provides information about concavity:

\[
f''(x)>0 \Rightarrow f \text{ is concave up},
\]
\[
f''(x)<0 \Rightarrow f \text{ is concave down}.
\]

Optimization uses derivatives to identify the maximum or minimum value of an objective function subject to constraints.

\subsubsection*{Learning Objectives}

Students will be able to:

\begin{enumerate}[label=\arabic*.]
    \item Find local and absolute extrema.
    \item Apply the Extreme Value Theorem and Mean Value Theorem.
    \item Determine intervals of increase and decrease.
    \item Determine concavity and locate possible inflection points.
    \item Sketch graphs using first- and second-derivative information.
    \item Solve optimization problems in biological, social, and economic contexts.
    \item Evaluate certain indeterminate limits using L'H\^{o}pital's Rule.
    \item Identify vertical, horizontal, and slant asymptotes.
    \item Approximate roots using Newton--Raphson iteration.
    \item Find antiderivatives of basic functions.
\end{enumerate}

\subsubsection*{Potentially Difficult Topics}

\begin{itemize}
    \item Translating optimization word problems into functions and constraints.
    \item Distinguishing local extrema from absolute extrema.
    \item Applying the first derivative test and second derivative test.
    \item Identifying true inflection points.
    \item Combining derivative information to sketch accurate graphs.
    \item Determining when L'H\^{o}pital's Rule is allowed.
    \item Newton--Raphson convergence and selecting suitable starting values.
\end{itemize}

\subsection{Module 5: Integral Calculus}

\textbf{Textbook basis:} Chapter 6.

\subsubsection*{Major Topics}

\begin{itemize}
    \item Area problem.
    \item Definite integrals.
    \item Riemann sums.
    \item Properties of the definite integral.
    \item Fundamental Theorem of Calculus.
    \item Antiderivatives and indefinite integrals.
    \item Cumulative change.
    \item Average value of a function.
    \item Area between curves.
    \item Volume of solids.
    \item Arc length and rectification of curves.
\end{itemize}

\subsubsection*{Core Concepts}

A definite integral represents accumulated change:
\[
\int_a^b f(x)\,dx.
\]

It can be approximated using Riemann sums:
\[
\int_a^b f(x)\,dx
=
\lim_{n\to \infty}
\sum_{i=1}^{n}f(x_i^*)\Delta x.
\]

The Fundamental Theorem of Calculus connects accumulation and rates of change. If:
\[
F'(x)=f(x),
\]
then:
\[
\int_a^b f(x)\,dx=F(b)-F(a).
\]

For an applied problem, if \(r(t)\) is a rate of change, then:
\[
\int_a^b r(t)\,dt
\]
represents the net accumulated change from time \(a\) to time \(b\).

\subsubsection*{Learning Objectives}

Students will be able to:

\begin{enumerate}[label=\arabic*.]
    \item Interpret the definite integral as net accumulation.
    \item Approximate integrals using Riemann sums.
    \item Use properties of definite integrals.
    \item Find antiderivatives.
    \item Apply both parts of the Fundamental Theorem of Calculus.
    \item Calculate total and net change from rate functions.
    \item Find average values of functions.
    \item Compute areas under curves and between curves.
    \item Interpret integration results with appropriate units.
\end{enumerate}

\subsubsection*{Potentially Difficult Topics}

\begin{itemize}
    \item Understanding the difference between signed area and total geometric area.
    \item Translating a Riemann sum into a definite integral.
    \item Applying the Fundamental Theorem of Calculus correctly.
    \item Choosing integration bounds.
    \item Interpreting integrals as accumulated quantities in context.
    \item Finding areas between curves when graphs intersect.
    \item Distinguishing definite and indefinite integrals.
\end{itemize}

\subsection{Module 6: Integration Techniques and Numerical Methods}

\textbf{Textbook basis:} Chapter 7.

\subsubsection*{Major Topics}

\begin{itemize}
    \item Substitution rule.
    \item Definite and indefinite integrals.
    \item Integration by parts.
    \item Rational functions.
    \item Partial-fraction decomposition.
    \item Improper integrals.
    \item Numerical integration.
    \item Midpoint rule.
    \item Trapezoidal rule.
    \item Error estimation.
    \item Taylor approximation.
    \item Tables of integrals.
\end{itemize}

\subsubsection*{Core Concepts}

Substitution reverses the chain rule. If:
\[
u=g(x),
\]
then:
\[
\int f(g(x))g'(x)\,dx=\int f(u)\,du.
\]

Integration by parts reverses the product rule:
\[
\int u\,dv=uv-\int v\,du.
\]

The trapezoidal rule approximates a definite integral numerically:
\[
\int_a^b f(x)\,dx
\approx
\frac{\Delta x}{2}
\left[
f(x_0)+2f(x_1)+\cdots+2f(x_{n-1})+f(x_n)
\right].
\]

\subsubsection*{Learning Objectives}

Students will be able to:

\begin{enumerate}[label=\arabic*.]
    \item Select an appropriate technique for evaluating an integral.
    \item Evaluate integrals using substitution.
    \item Use integration by parts.
    \item Decompose rational functions into partial fractions.
    \item Evaluate selected improper integrals.
    \item Approximate definite integrals numerically.
    \item Compare numerical approximations with exact values when available.
    \item Interpret numerical-integration results in applied settings.
    \item Use approximation methods when closed-form antiderivatives are unavailable.
\end{enumerate}

\subsubsection*{Potentially Difficult Topics}

\begin{itemize}
    \item Recognizing an appropriate substitution.
    \item Changing bounds during definite-integral substitution.
    \item Choosing \(u\) and \(dv\) for integration by parts.
    \item Partial-fraction decomposition.
    \item Understanding convergence of improper integrals.
    \item Applying midpoint and trapezoidal rules accurately.
    \item Estimating and interpreting numerical error.
\end{itemize}

\subsection{Module 7: Introduction to Differential Equations and Applied Modeling}

\textbf{Textbook basis:} Selected material from Chapter 8.

\subsubsection*{Major Topics}

\begin{itemize}
    \item Separable differential equations.
    \item Pure-time differential equations.
    \item Autonomous differential equations.
    \item Equilibrium points.
    \item Stability of equilibria.
    \item Direction fields and vector fields.
    \item Compartment models.
    \item Ecological models.
    \item Chemical reaction models.
    \item Epidemic models.
    \item Integrating factors.
    \item Two-compartment models.
\end{itemize}

\subsubsection*{Core Concepts}

A differential equation relates a quantity to its rate of change:
\[
\frac{dy}{dt}=f(t,y).
\]

A separable equation can be written in the form:
\[
\frac{dy}{dt}=g(t)h(y),
\]
which allows separation:
\[
\frac{1}{h(y)}\,dy=g(t)\,dt.
\]

An equilibrium is a value \(y^*\) satisfying:
\[
f(y^*)=0.
\]

Equilibria can be stable or unstable. A stable equilibrium attracts nearby solutions, while an unstable equilibrium repels them.

\subsubsection*{Learning Objectives}

Students will be able to:

\begin{enumerate}[label=\arabic*.]
    \item Interpret differential equations as models of change.
    \item Solve elementary separable differential equations.
    \item Apply initial conditions.
    \item Identify equilibrium solutions.
    \item Determine stability of equilibria using signs and graphs.
    \item Interpret direction fields.
    \item Construct simple population, compartment, and epidemic models.
    \item Explain model parameters in scientific or social-science contexts.
\end{enumerate}

\subsubsection*{Potentially Difficult Topics}

\begin{itemize}
    \item Separating variables correctly.
    \item Applying initial conditions after integration.
    \item Distinguishing equilibrium solutions from general solutions.
    \item Interpreting stability graphically.
    \item Connecting equations, slope fields, and real-world behavior.
    \item Modeling assumptions and limitations.
\end{itemize}

\section{Suggested Weekly Schedule}

\begin{longtable}{p{0.1\textwidth} p{0.28\textwidth} p{0.52\textwidth}}
\toprule
\textbf{Week} & \textbf{Module / Chapter} & \textbf{Topics and Expected Focus} \\
\midrule
\endfirsthead

\toprule
\textbf{Week} & \textbf{Module / Chapter} & \textbf{Topics and Expected Focus} \\
\midrule
\endhead

1 & Module 0: Chapter 1 & Precalculus diagnostic; algebra review; functions; function notation; domain and range. \\
2 & Module 0: Chapter 1 & Graphs, transformations, polynomial and rational functions, exponentials, logarithms, and logarithmic scales. \\
3 & Module 1: Chapter 2 & Discrete-time models, exponential growth and decay, recurrence equations, sequences. \\
4 & Module 1: Chapter 2 & Density dependence, discrete logistic models, applied population and drug-absorption examples. \\
5 & Module 2: Chapter 3 & Limits, graphical and numerical limits, algebraic limit evaluation, one-sided limits. \\
6 & Module 2: Chapter 3 & Continuity, limits at infinity, asymptotic behavior, Intermediate Value Theorem. \\
7 & Module 3: Chapter 4 & Derivative definition, derivative interpretation, power rule, polynomial and rational derivatives. \\
8 & Module 3: Chapter 4 & Product rule, quotient rule, chain rule, higher derivatives. \\
9 & Module 3: Chapter 4 & Exponential, logarithmic, and trigonometric derivatives; implicit differentiation; related rates. \\
10 & Module 4: Chapter 5 & Extrema, Mean Value Theorem, monotonicity, concavity, and graphing. \\
11 & Module 4: Chapter 5 & Optimization, L'H\^{o}pital's Rule, linear approximation, Newton--Raphson method. \\
12 & Module 5: Chapter 6 & Definite integrals, Riemann sums, accumulation, and properties of integrals. \\
13 & Module 5: Chapter 6 & Fundamental Theorem of Calculus, antiderivatives, cumulative change, average values. \\
14 & Module 6: Chapter 7 & Substitution, integration by parts, numerical integration, and applied integral problems. \\
15 & Module 7: Chapter 8 & Introductory separable differential equations, equilibrium, stability, and modeling review. \\
\bottomrule
\end{longtable}

\section{Concept Dependency Map}

\[
\text{Algebra and Functions}
\longrightarrow
\text{Limits}
\longrightarrow
\text{Derivatives}
\longrightarrow
\text{Applications of Derivatives}
\longrightarrow
\text{Antiderivatives}
\longrightarrow
\text{Definite Integrals}
\longrightarrow
\text{Differential Equations and Modeling}.
\]

A more detailed dependency sequence is:

\begin{enumerate}[label=\arabic*.]
    \item Algebra, graphs, and elementary functions.
    \item Exponential and logarithmic functions.
    \item Limits and continuity.
    \item Derivative definition and differentiation rules.
    \item Chain rule and implicit differentiation.
    \item Derivative applications: optimization, rates, graphing, approximation.
    \item Antiderivatives and accumulation.
    \item Definite integrals and the Fundamental Theorem of Calculus.
    \item Integration techniques and numerical methods.
    \item Differential-equation models and stability.
\end{enumerate}

\section{Most Challenging Topics and Support Strategy}

\begin{longtable}{p{0.28\textwidth} p{0.3\textwidth} p{0.32\textwidth}}
\toprule
\textbf{Topic} & \textbf{Why It Is Difficult} & \textbf{Recommended Support} \\
\midrule
\endfirsthead

\toprule
\textbf{Topic} & \textbf{Why It Is Difficult} & \textbf{Recommended Support} \\
\midrule
\endhead

Limits & Requires understanding behavior near a point rather than only substitution. & Use graphs, numerical tables, and algebraic examples together. \\
Chain Rule & Students must recognize nested functions and account for inner-function change. & Practice identifying outer and inner functions before differentiating. \\
Implicit Differentiation & Variables are mixed together and derivatives of \(y\) require careful notation. & Differentiate term by term and explicitly write \(dy/dx\). \\
Related Rates & Requires translating a physical or biological situation into variables and equations. & Draw diagrams, define units, and differentiate with respect to time. \\
Optimization & Requires modeling, constraints, derivatives, and interpretation of answers. & Use a structured process: variables, objective, constraint, substitution, derivative, interpretation. \\
Riemann Sums & Connects finite sums with the abstract idea of an integral. & Begin with rectangles and graphical area estimates. \\
Fundamental Theorem of Calculus & Connects derivatives and integrals, which students often view as unrelated operations. & Emphasize accumulation functions and rate-versus-total-change examples. \\
Substitution & Requires recognizing reverse chain-rule structure. & Practice matching \(g(x)\) with \(g'(x)\). \\
Differential Equations & Requires interpreting solutions and parameters, not merely computing. & Use population, drug, and epidemic examples with graphs. \\
\bottomrule
\end{longtable}

\section{Applications for Biological and Social Sciences}

\subsection{Biological Science Applications}

\begin{itemize}
    \item Population growth and carrying capacity.
    \item Predator--prey and ecological interactions.
    \item Drug absorption and passage of drugs through the human body.
    \item Chemical reaction rates.
    \item Epidemic and disease-spread models.
    \item Diffusion and movement of substances.
    \item Compartment models for nutrients, drugs, or organisms.
    \item Microbial growth and enzyme activity.
\end{itemize}

\subsection{Social Science Applications}

\begin{itemize}
    \item Population and demographic change.
    \item Exponential growth and decay in social indicators.
    \item Marginal analysis in economics.
    \item Optimization of cost, revenue, utility, and resource allocation.
    \item Modeling cumulative quantities over time.
    \item Rates of change in income, employment, production, or demand.
    \item Linearization and approximation for interpreting small changes.
    \item Numerical methods for data-based decisions.
\end{itemize}

\section{Assessment Suggestions}

\begin{table}[h!]
\centering
\begin{tabular}{p{0.42\textwidth} p{0.2\textwidth} p{0.28\textwidth}}
\toprule
\textbf{Assessment Component} & \textbf{Suggested Weight} & \textbf{Purpose} \\
\midrule
Homework and online practice & 15--20\% & Build procedural fluency and reinforce concepts. \\
Quizzes & 10--15\% & Check weekly understanding and identify misconceptions early. \\
Modeling assignments or labs & 10--15\% & Apply calculus to biological and social-science contexts. \\
Midterm examinations & 20--30\% & Assess foundational and derivative-based learning. \\
Final examination & 25--35\% & Assess cumulative mastery of limits, derivatives, and integrals. \\
Participation, worksheets, or projects & 5--10\% & Encourage discussion, interpretation, and collaborative problem solving. \\
\bottomrule
\end{tabular}
\end{table}

\section{Recommended Pedagogical Emphases}

\begin{itemize}
    \item Use real data and contextual problems whenever possible.
    \item Connect symbolic work to graphs, tables, and verbal interpretation.
    \item Emphasize units for derivatives and integrals.
    \item Distinguish clearly between discrete-time and continuous-time models.
    \item Require students to explain the meaning of answers, not only calculate them.
    \item Use technology for graphing, numerical approximation, recurrence equations, and data visualization.
    \item Include population, epidemiological, pharmacokinetic, ecological, and economic examples.
    \item Provide frequent cumulative review because derivative and integration skills build sequentially.
\end{itemize}

\section{Optional Extension Topics}

The following textbook topics are valuable but are more appropriate as extensions, enrichment material, or content for later courses:

\begin{itemize}
    \item Formal \(\varepsilon\)--\(\delta\) limit definitions.
    \item Rigorous theory of Riemann integration.
    \item Improper integrals.
    \item Taylor polynomials and approximation error.
    \item Systems of recurrence equations.
    \item Linear algebra, matrices, eigenvalues, and eigenvectors.
    \item Multivariable calculus.
    \item Systems of differential equations.
    \item Probability and statistics.
\end{itemize}

\section{Summary}

The recommended Calculus I course should concentrate on Chapters 1--7 of the textbook, with selected introductory material from Chapter 8. The course progresses from functions and modeling to limits, derivatives, applications of derivatives, integration, integration techniques, and basic differential equations.

The central organizing ideas are:

\[
\text{Functions}
\rightarrow
\text{Change}
\rightarrow
\text{Rates of Change}
\rightarrow
\text{Accumulation}
\rightarrow
\text{Mathematical Modeling}.
\]

For biological and social science students, the course should prioritize conceptual understanding, applied interpretation, model construction, graphical reasoning, and communication of quantitative results.

\end{document}